\chapter{Parser streaming + batching}

\section{Fluxo desta etapa}

Agora juntamos streaming, dataclass e Postgres.

\begin{verbatim}
XML element
    -> parse attributes
    -> HealthRecord dataclass
    -> batch list
    -> Postgres
\end{verbatim}

\section{Parser de elemento Record}

Arquivo: \code{src/parser/xml_parser.py}

\begin{lstlisting}[style=devbutterpython]
import xml.etree.ElementTree as ET
from src.models.health_record import HealthRecord


def clean_tag(tag: str) -> str:
    if "}" in tag:
        return tag.split("}", 1)[1]
    return tag


def safe_float(value: str | None) -> float | None:
    if value is None:
        return None

    try:
        return float(value)
    except ValueError:
        return None


def extract_metadata(element: ET.Element) -> dict[str, str]:
    metadata: dict[str, str] = {}

    for child in element:
        if clean_tag(child.tag) != "MetadataEntry":
            continue

        key = child.attrib.get("key")
        value = child.attrib.get("value")

        if key is not None and value is not None:
            metadata[key] = value

    return metadata


def parse_record_element(
    element: ET.Element,
    record_id: int,
) -> HealthRecord:
    raw_value = element.attrib.get("value")

    return HealthRecord(
        apple_record_id=record_id,
        type=element.attrib.get("type"),
        source_name=element.attrib.get("sourceName"),
        source_version=element.attrib.get("sourceVersion"),
        device=element.attrib.get("device"),
        unit=element.attrib.get("unit"),
        creation_date=element.attrib.get("creationDate"),
        start_date=element.attrib.get("startDate"),
        end_date=element.attrib.get("endDate"),
        value=raw_value,
        value_numeric=safe_float(raw_value),
        metadata=extract_metadata(element),
    )
\end{lstlisting}

\section{Repository com batch insert}

Instale o driver:

\begin{lstlisting}[style=devbutterbash]
poetry add "psycopg[binary]"
\end{lstlisting}

Arquivo: \code{src/db/repository.py}

\begin{lstlisting}[style=devbutterpython]
import json
from collections.abc import Sequence

import psycopg

from src.models.health_record import HealthRecord


class HealthRecordRepository:
    def __init__(self, connection: psycopg.Connection) -> None:
        self.connection = connection

    def insert_batch(self, records: Sequence[HealthRecord]) -> None:
        if not records:
            return

        rows = [
            (
                record.apple_record_id,
                record.type,
                record.source_name,
                record.source_version,
                record.device,
                record.unit,
                record.creation_date,
                record.start_date,
                record.end_date,
                record.value,
                record.value_numeric,
                json.dumps(record.metadata),
            )
            for record in records
        ]

        with self.connection.cursor() as cursor:
            cursor.executemany(
                """
                INSERT INTO health_records (
                    apple_record_id,
                    type,
                    source_name,
                    source_version,
                    device,
                    unit,
                    creation_date,
                    start_date,
                    end_date,
                    value,
                    value_numeric,
                    metadata
                )
                VALUES (
                    %s, %s, %s, %s, %s, %s,
                    %s, %s, %s, %s, %s, %s
                )
                """,
                rows,
            )

        self.connection.commit()
\end{lstlisting}

\section{Pipeline sincrona}

Arquivo: \code{src/pipelines/sync_pipeline.py}

\begin{lstlisting}[style=devbutterpython]
from pathlib import Path
from time import perf_counter
import logging
import xml.etree.ElementTree as ET

from src.db.repository import HealthRecordRepository
from src.parser.xml_parser import clean_tag, parse_record_element

logger = logging.getLogger(__name__)


def run_sync_pipeline(
    xml_file: Path,
    repository: HealthRecordRepository,
    batch_size: int,
) -> None:
    batch = []
    record_id = 0
    started_at = perf_counter()

    context = ET.iterparse(xml_file, events=("end",))

    for _, element in context:
        tag = clean_tag(element.tag)

        if tag == "Record":
            record_id += 1
            batch.append(parse_record_element(element, record_id))

            if len(batch) >= batch_size:
                insert_started_at = perf_counter()
                repository.insert_batch(batch)
                insert_elapsed = perf_counter() - insert_started_at

                logger.info(
                    "Inserted %s records in %.2fs | total=%s",
                    len(batch),
                    insert_elapsed,
                    record_id,
                )

                batch.clear()

        element.clear()

    if batch:
        repository.insert_batch(batch)

    elapsed = perf_counter() - started_at
    logger.info(
        "Finished parsing %s records in %.2fs | %.2f records/s",
        record_id,
        elapsed,
        record_id / elapsed if elapsed > 0 else 0,
    )
\end{lstlisting}

\section{Conceitos aplicados}

\begin{itemize}
  \item \code{batch}: uma \code{list}; append e O(1) medio.
  \item \code{batch_size}: tradeoff entre memoria e velocidade.
  \item \code{iterparse}: lazy loading / streaming.
  \item \code{insert_batch}: reduz overhead de I/O com o banco.
  \item \code{perf_counter}: mede tempo com precisao suficiente para benchmark.
\end{itemize}

\section{Experimento}

Rode a pipeline com:

\begin{verbatim}
batch_size = 500
batch_size = 5000
batch_size = 20000
batch_size = 50000
\end{verbatim}

Anote:

\begin{itemize}
  \item tempo total;
  \item records/s;
  \item memoria percebida;
  \item tempo medio por batch.
\end{itemize}
